\scanpage{151}
\[
A(T)=\mathbb Z(\widehat T)
\]

Soit $\rho\unsure{\in|:}T\longrightarrow\operatorname{Gl}(n)$ / induit $\rho:T\longrightarrow k^{*\,n}$, d'où $n$ éléments
$\alpha_1,\ldots\alpha_n$ de $\widehat T$. Les fonctions symétriques élémentaires en
les $\alpha_i$ sont notées $C^i(\rho)$ $(0\leq i\leq n)$, ils ne dépendent
que de la classe de $\rho$. Soit $\unsure{C|\mathcal C}(\rho)=\sum_0^\infty C^i(\rho)\in\mathbb Z[\widehat T]$
($C^i(\unsure{\varphi|\rho})=0$ si $i>\deg\rho$) ; si $\rho$ est une extension de $\rho'$ par
$\rho''$, on a $C(\rho)=C(\rho')C(\rho'')$.

On peut prolonger $\rho\mapsto(\deg\rho,C(\rho))$ en un hom de groupe
\[
A(T)=\mathbb Z(\widehat T)\longrightarrow\mathbb Z\times(1+\mathbb Z[[\widehat T]]^+) \tag*{(1)}
\]
c'est un hom d'anneaux de Chern, caractérisé par
ses valeurs sur les $\rho_\chi$ $(\chi\in\widehat T)$
\[
\widetilde C(\rho_\chi)=(1,1+[\chi])
\]
Je dis que cet hom (1) est injectif \deleted{\unsure{(1)}} et à image
\deleted{\unsure{discr.}}dense. Pour le voir, on compose avec \deleted{\unsure{$\varphi$}}
\[
\varphi:\mathbb Z\times(1+\mathbb Z[[\widehat T]]^+)\longrightarrow\mathbb Q[[\widehat T]]
\]
transformant $\chi$ en $e^\chi$, et on note que les
$e^\chi$ sont linéairement indépendants. Pour montrer que
l'image \deleted{est} dense, on montre qu'on peut trouver
un hom $\rho$ dont les $C^1,\ldots C^n\unsure{\mid|)}$ soient donnés,
i.e \deleted{des formes} étant donné des \deleted{formes sur $\mathbb Z^r$}
polynomes $C^i\in\mathbb Z[X_1,\ldots,X_r]$ $(1\leq i\leq n)$, $C^i$
de degré $i$, on peut trouver des formes linéaires
$\alpha^1,\ldots,\alpha^N$ en les $X_i$, telles que
$C^i=\sigma_i(\alpha^1,\ldots,\alpha^N)$ pour $1\leq i\leq n$. Raisonnons sur
% archive sheet number: 150
% printed form (vertical paper watermark, reverse face)
\tikz[baseline=-.4ex,scale=.15]{\foreach \a in {0,45,...,315}{\draw (\a:.52) circle (.22);} \draw (0,0) circle (.21);} ARGENTINE \tikz[baseline=-.2ex,scale=.17]{\draw (0,0)--(1.6,0)--(1.8,.33)--(.5,.33)--(.32,.13)--(.92,.13);} TYPEWRITER \tikz[baseline=-.4ex,scale=.15]{\foreach \a in {0,45,...,315}{\draw (\a:.52) circle (.22);} \draw (0,0) circle (.21);} \unsure{AR}
